
\section{Introducción}
En esta actividad aprendimos a usar y utilizamos los ''Issues'', estos mismos son herramientas de GitHub nos permite realizar varias acciones, cuales como reportar errores que puede generar un código mal implementado por nuestros colaboradores, también nos permite agregar tareas nuevas a los repositorios en donde nos encontremos trabajando, dentro de un ambiente colaborativo, al igual que podemos realizar la petición de la mejora de algún código incompleto o que puede llegar a complementarse a un mejor proyecto de software, otra de sus funciones más representativas es la posibilidad de crear una buena organización en un equipo para la realización de un trabajo colaborativo con el fin de realizar un proyecto.
\\En esta practica se realizó la creación de un Issue desde dos maneras posibles la primera fue la creación un Issue por medio de una pagina web y la segunda forma fue crear un Issue desde la terminal de VS code, las dos maneras de realizar esta actividad tienen su propia forma de realización, sus comandos aplicables, cada miembro del equipo participo en la creación de los Issues en las dos formas descritas en el documento de la actividad propuesta.
